% !TEX program = lualatex
% Transparencias de Fundamentos de las Matemáticas II
% Clase 01: Números enteros I

\documentclass[aspectratio=169,11pt]{beamer}

\usepackage{fontspec}
\usepackage[spanish,es-nodecimaldot]{babel}
\usepackage{amsmath,amssymb,mathtools}
\usepackage{booktabs}
\usepackage{csquotes}
\usepackage{xcolor}

\usetheme{Madrid}
\usecolortheme{default}
\setbeamertemplate{navigation symbols}{}
\setbeamertemplate{footline}[frame number]

\definecolor{FdMBlue}{RGB}{20,55,100}
\definecolor{FdMRed}{RGB}{130,30,30}
\definecolor{FdMGreen}{RGB}{25,100,60}

\setbeamercolor{structure}{fg=FdMBlue}
\setbeamercolor{title}{fg=white,bg=FdMBlue}
\setbeamercolor{frametitle}{fg=white,bg=FdMBlue}
\setbeamercolor{block title}{fg=white,bg=FdMBlue}
\setbeamercolor{block body}{bg=blue!3}
\setbeamercolor{alerted text}{fg=FdMRed}

\setmainfont{Latin Modern Roman}
\setsansfont{Latin Modern Sans}
\setmonofont{Latin Modern Mono}

\newcommand{\N}{\mathbb{N}}
\newcommand{\Z}{\mathbb{Z}}
\newcommand{\Q}{\mathbb{Q}}
\newcommand{\R}{\mathbb{R}}
\newcommand{\C}{\mathbb{C}}
\newcommand{\Pow}{\mathcal{P}}
\newcommand{\id}{\operatorname{id}}
\newcommand{\mcd}{\operatorname{mcd}}
\newcommand{\mcm}{\operatorname{mcm}}
\newcommand{\im}{\operatorname{Im}}

\title[FdM II -- Clase 01]{Números enteros I}
\subtitle{Ampliación de los naturales, opuesto y orden}
\author{Antonio Falcó Montesinos}
\institute{Universidad CEU Cardenal Herrera}
\date{Fundamentos de las Matemáticas II}

\begin{document}

\begin{frame}
  \titlepage
\end{frame}

\begin{frame}{Objetivos de la clase}
\begin{itemize}
  \item Motivar la construcción de \(\mathbb{Z}\).
  \item Distinguir naturales, enteros positivos, cero y enteros negativos.
  \item Trabajar con opuesto, suma y orden en \(\mathbb{Z}\).
\end{itemize}
\end{frame}

\begin{frame}{Mapa de la clase}
\tableofcontents
\end{frame}

\section{Motivación}

\begin{frame}{Motivación}
\begin{block}{Por qué ampliar \(\mathbb{N}\)}
\begin{itemize}
  \item En \(\mathbb{N}\), ecuaciones como \(10+x=7\) no tienen solución.
  \item Los enteros introducen números que representan diferencias negativas.
  \item La ampliación responde a una necesidad algebraica concreta.
\end{itemize}
\end{block}
\end{frame}

\begin{frame}{Motivación}
\begin{block}{El conjunto \(\mathbb{Z}\)}
\begin{itemize}
  \item \(\mathbb{Z}=\{\ldots,-3,-2,-1,0,1,2,3,\ldots\}\).
  \item Todo natural es entero: \(\mathbb{N}\subset\mathbb{Z}\).
  \item No todo entero es natural: \(-4\in\mathbb{Z}\), pero \(-4\notin\mathbb{N}\).
\end{itemize}
\end{block}
\end{frame}

\begin{frame}{Motivación}
\begin{block}{Opuesto}
\begin{itemize}
  \item Cada entero \(a\) tiene un opuesto \(-a\).
  \item Se cumple \(a+(-a)=0\).
  \item El cero es su propio opuesto: \(-0=0\).
\end{itemize}
\end{block}
\end{frame}

\section{Orden y operaciones}

\begin{frame}{Orden y operaciones}
\begin{block}{Suma y producto}
\begin{itemize}
  \item Si \(a,b\in\mathbb{Z}\), entonces \(a+b\in\mathbb{Z}\) y \(ab\in\mathbb{Z}\).
  \item La suma y el producto son operaciones internas en \(\mathbb{Z}\).
  \item Estas operaciones extienden las operaciones conocidas en \(\mathbb{N}\).
\end{itemize}
\end{block}
\end{frame}

\begin{frame}{Orden y operaciones}
\begin{block}{Orden}
\begin{itemize}
  \item El orden en \(\mathbb{Z}\) permite comparar enteros positivos y negativos.
  \item Si \(a<b\), entonces \(a+c<b+c\).
  \item Multiplicar por un entero negativo invierte desigualdades.
\end{itemize}
\end{block}
\end{frame}

\begin{frame}{Orden y operaciones}
\begin{block}{Errores frecuentes}
\begin{itemize}
  \item Interpretar \(-3\) como una resta pendiente.
  \item Olvidar que \(0\) no es positivo ni negativo.
  \item Suponer que las reglas de signos son convenciones sin justificación estructural.
\end{itemize}
\end{block}
\end{frame}

\section{Profundización}

\begin{frame}{Profundización}
\begin{block}{Modelo de lectura}
\begin{itemize}
  \item Antes de usar \(a< b\), declarar que \(a,b\in\mathbb{Z}\).
  \item Antes de multiplicar desigualdades, revisar el signo del factor.
  \item Distinguir propiedades de suma y propiedades de producto.
\end{itemize}
\end{block}
\end{frame}

\begin{frame}{Profundización}
\begin{block}{Ejemplo guiado}
\begin{itemize}
  \item Resolver \(5+x=-2\) en \(\mathbb{Z}\).
  \item Sumamos el opuesto de \(5\): \(x=-2-5\).
  \item Por tanto, \(x=-7\).
\end{itemize}
\end{block}
\end{frame}

\begin{frame}{Profundización}
\begin{block}{Pregunta de control}
\begin{itemize}
  \item ¿Qué ecuaciones tienen solución en \(\mathbb{N}\) pero no en \(\mathbb{Z}\)?
  \item ¿Qué ecuaciones no tienen solución en \(\mathbb{N}\) y sí en \(\mathbb{Z}\)?
  \item ¿Qué propiedad algebraica aporta el opuesto?
\end{itemize}
\end{block}
\end{frame}


\section{Detalle y práctica}

\begin{frame}{Detalle y práctica}
\begin{block}{Construcción por diferencias}
\begin{itemize}
  \item Una forma rigurosa de construir \(\mathbb{Z}\) consiste en representar un entero como una diferencia \(a-b\), con \(a,b\in\mathbb{N}\).
  \item Pares distintos pueden representar el mismo entero: \(5-2\) y \(7-4\) representan \(3\).
  \item La igualdad correcta es \(a-b=c-d\) si y solo si \(a+d=b+c\).
\end{itemize}
\end{block}
\end{frame}

\begin{frame}{Detalle y práctica}
\begin{block}{Actividad breve}
\begin{itemize}
  \item Representar \(-3\) mediante tres diferencias distintas de naturales.
  \item Decidir si los pares \((2,7)\) y \((5,10)\) representan el mismo entero.
  \item Justificar por qué la suma de clases de diferencias está bien definida.
\end{itemize}
\end{block}
\end{frame}

\begin{frame}{Cierre}
\begin{itemize}
  \item Los enteros permiten resolver ecuaciones aditivas que no se resuelven en \(\mathbb{N}\).
  \item La suma y el producto siguen siendo operaciones internas.
  \item El orden exige cuidado al multiplicar por números negativos.
\end{itemize}
\end{frame}

\begin{frame}{Trabajo recomendado}
\begin{itemize}
  \item Releer las secciones correspondientes del manual.
  \item Identificar definiciones, símbolos nuevos y enunciados principales.
  \item Preparar dudas y ejemplos para el seminario asociado.
\end{itemize}
\end{frame}

\end{document}
